\section{Excluding finite abelian characters}\label{sec:characters}

Throughout this section, $D$ is a central division algebra of odd degree
$n\geq3$ over the quadratic extension $L/k$, with unitary involution $*$.
We write $U=\U(D,*)$ and $S=\SU(D,*)$. The remaining alternative in
Proposition~\ref{quo:dichotomy} is excluded by the following result.

\begin{theorem}\label{char:trivial}
Every homomorphism from $U(k)$ to a finite abelian group is trivial on
$S(k)$.
\end{theorem}

The difference-function argument below adapts the corresponding argument
of \cite[Section~4]{MPNF}. Its two arithmetic ingredients have particularly
useful function-field forms: every nonsingular Hermitian element admits
a scalar normalization, and the additive span of $U(k)$ is all of $D$.
The latter is an equality of abstract additive groups. Local density alone
would not suffice for the exact transformations used here.
Fix a homomorphism $\chi:U(k)\to B$, with $B$ finite and written additively.
We construct a function on $D$ whose differences are controlled by
$\chi$, and prove that exact transformations of pairs make those
differences constant. Scalar normalization then gives a function on
all of $D^\times$, invariant under a noncentral normal subgroup of
the inner group $\SL_1(D)(L)$. The known inner-type theorem turns this
invariance into continuity of $\chi|_{S(k)}$ at $A$; local unitary
commutators then force its vanishing.

\subsection{Hermitian normalization and a cyclic maximal field}

We first arrange a commutative field inside $D$ on which both conjugation
and the involution can be calculated explicitly.

\begin{lemma}\label{char:normalization}
If $h=h^*\ne0$ and $s=\Nrd(h)\in k^\times$, there exists
$a\in D^\times$ such that
\begin{equation}\label{char:normalized-norm}
 a^*a=sh,\qquad \Nrd(a)=s^{(n+1)/2}.
\end{equation}
There is also a $*$-stable maximal subfield $E\subset D$ of the form
$E=LF$, where $F/k$ is cyclic of degree $n$ and $*$ restricts to the
nonidentity automorphism of $E/F$.
\end{lemma}

\begin{proof}
The variety defined by Equation~\ref{char:normalized-norm} is a torsor
under $S$: over a separable closure the Hermitian equation has a solution,
and the prescribed determinant is compatible because
\[
 \Nrd(sh)=s^{n+1}=N_{L/k}\bigl(s^{(n+1)/2}\bigr).
\]
The simply connected $H^1$-vanishing theorem over global function fields
therefore gives a rational point \cite{Harder,Conrad}. This applies without
any restriction on the characteristic or on its divisibility into $n$.

The cyclic approximation theorem for global fields supplies a cyclic
extension $F/k$ of degree $n$ with full local degree at each place
below the Brauer support of $D$
\cite[Theorems~1--2 and Section~4]{LorenzRoquette}. Prescribing a
degree-$n$ local field at one of these places ensures that the global
cyclic Galois algebra is a field. The characteristic-primary part is covered
by Artin--Schreier--Witt approximation, and the Wang condition is
automatic in positive characteristic. Since $n$ is odd, $F$ and $L$
are disjoint.
The local invariant criterion shows that $LF$ splits $D$. Its degree
over $L$ is $n$, so the embedding criterion gives an embedding
$\alpha:LF\hookrightarrow D$ \cite{Reiner}.

Let bar on $LF$ fix $F$ and induce the nonidentity automorphism of $L/k$.
Skolem--Noether provides $h\in D^\times$ with
\[
 \alpha(\bar z)^*=h\alpha(z)h^{-1}\qquad(z\in LF).
\]
Taking adjoints shows that $h^*$ is another intertwiner, whence
$c=h^{-1}h^*\in\alpha(LF)^\times$. On this field
$x^*=h\bar xh^{-1}$, and consequently $\bar c=c^{-1}$.
Hilbert 90 gives $t\in\alpha(LF)^\times$ with $t/\bar t=c$.
Replacing $h$ by $ht$ makes it Hermitian, since
$(ht)^*=h\bar t c=ht$, and preserves the intertwining identity.
Choose $a$ as in Equation~\ref{char:normalized-norm} for this $h$.
Then the embedding $z\mapsto a\alpha(z)a^{-1}$ is $*$-stable:
\[
 (a\alpha(\bar z)a^{-1})^*
 =(a^*)^{-1}h\alpha(z)h^{-1}a^*
 =a\alpha(z)a^{-1}.
\]
Its image is the required field $E$.
\end{proof}

Write $\bar z=z^*$ on $E$, and put
$E^1=\{z\in E^\times:z\bar z=1\}$. This group has infinitely many
points, by the Hilbert 90 parametrization. We shall use parameters
$d\in E^\times$ with the following properties:
\begin{equation}\label{char:parameter-conditions}
\begin{gathered}
 t_1,\ldots,t_n\text{, the $L$-conjugates of $d$, are distinct},\\
 \{t_1,\ldots,t_n\}\cap\{\bar t_1,\ldots,\bar t_n\}=\varnothing,
 \qquad d\bar d\notin k.
\end{gathered}
\end{equation}
The last condition implies $t_j\bar t_j\ne1$ for every $j$.
It ensures separability of a degree-two parameter map even in
characteristic two.

\subsection{A difference function and its exact identities}

Let $\mathcal D$ be the set of nonzero $b\in D$ such that
$b+b^*=a^*a$ for some $a\in D$. Define
\begin{equation}\label{char:f}
 m_b=1-ab^{-1}a^*,\qquad f(b)=\chi(m_b).
\end{equation}
Multiplication gives $m_bm_b^*=1$. If $a'{}^*a'=a^*a\ne0$, then
$a'=va$ for $v\in U(k)$, so the resulting elements $m_b$ are conjugate.
When $a=0$, $m_b=1$. Thus $f$ is well-defined, including on skew
elements. Using $ax$ in the defining equation also gives
\begin{equation}\label{char:congruence}
 f(x^*bx)=f(b)\qquad(x\in D^\times).
\end{equation}

Fix $d$ satisfying Equation~\ref{char:parameter-conditions}.
For $(a,c)\in D^2\setminus\{(0,0)\}$ there is a unique $e\in D$ with
\begin{equation}\label{char:sylvester}
 e+e^*=c^*c,\qquad ed+d^*e^*=a^*a.
\end{equation}
Indeed, solve the linear equation
\[
 ed-d^*e=a^*a-d^*c^*c.
\]
After splitting $D$, its linear part is a Sylvester operator whose
eigenvalues are the differences of the disjoint reduced spectra of
$d$ and $d^*$. It is therefore invertible. Taking adjoints shows that
$c^*c-e^*$ is a second solution of the same linear equation, proving
both identities in Equation~\ref{char:sylvester}. The solution is
nonzero, since $e=0$ would force $a=c=0$. We may consequently define
\[
 F_d(a,c)=f(ed)-f(e).
\]

\begin{lemma}\label{char:pair-invariance}
The function $F_d$ is invariant under independent left multiplications
of its two arguments by elements of $U(k)$ and under the transformation
\begin{equation}\label{char:pair-move}
 (a,c)\longmapsto(a-cd,a-c).
\end{equation}
\end{lemma}

\begin{proof}
The unitary multiplications leave Equation~\ref{char:sylvester}
unchanged. For the pair transformation put $b=ed$ and
$z=b^*+e-a^*c$. Expanding in the indicated order gives
\[
 z+z^*=(a-c)^*(a-c),\qquad
 zd+d^*z^*=(a-cd)^*(a-cd).
\]
The pair transformation is invertible, since $d\ne1$, so $z\ne0$.
Set $P=ce^{-1}-ab^{-1}$ and $Q=b^{-*}ze^{-1}$, where
$b^{-*}=(b^*)^{-1}$. Direct multiplication gives
\[
 P^*P=Q+Q^*,\qquad Q=(b^{-1})^*(zd)b^{-1}.
\]
The multiplicative identity that relates their characters is
\begin{equation}\label{char:three-unitaries}
 1-PQ^{-1}P^*=m_bm_zm_e^*,
\end{equation}
where the defining vectors for $m_b,m_z,m_e$ are $a,c-a,c$,
respectively. One can check the identity without commuting any factors:
use $Q^{-1}=ez^{-1}b^*$ and
\begin{align*}
 Pe&=-ab^{-1}z+m_b(c-a),\\
 b^*P^*&=ze^{-*}c^*+(c-a)^*m_e^*.
\end{align*}
In multiplying these expressions, the terms not containing
$-m_b(c-a)z^{-1}(c-a)^*m_e^*$ sum to $m_bm_e^*$; here one uses
$m_ba=-ab^{-1}b^*$ and $c^*c=e+e^*$.
By congruence invariance the left side of
Equation~\ref{char:three-unitaries} has character $f(zd)$.
Hence $f(zd)=f(b)+f(z)-f(e)$, which is precisely the assertion.
\end{proof}

To make this invariance useful, we must show that the transformations
act transitively on nonzero pairs. We first prove the additive assertion
that will turn one nonzero elementary shear into arbitrary addition.

\begin{lemma}\label{char:additive-span}
The additive subgroup generated by $U(k)$ equals $D$:
\[
 \bbZ U(k)=D.
\]
\end{lemma}

\begin{proof}
Put $M=\bbZ U(k)$; it is a subring because $U(k)$ is a group. Choose
a place of $k$ split in $L$. Hilbert 90 and weak approximation give
a central phase $z\in L^1$ with a simple pole at one of the two places
above it and a simple zero at the other. Then
$t_0=z+z^{-1}\in k\cap M$ has a simple pole at the chosen place.
In particular $k/ k_0$ is finite separable for
$k_0=\bbF_p(t_0)$, where $p=\operatorname{char}(k)$.
The group $M$ is a module over $R=\bbF_p[t_0]$.

Fix a finite prime $\pi$ of $R$, let $K$ be the corresponding
completion of $k_0$, and let $C$ be the additive closure of $M$ in
$D_K=D\otimes_{k_0}K$. It is a closed module over the valuation ring
$\mathcal O_K$. Weak approximation for $U$ puts the product of all
its local groups above $\pi$ inside $C$. The largest $K$-vector
subspace contained in $C$ is
\[
 W=\bigcap_{r\geq0}\pi^r C.
\]
Indeed the right side is closed, is an $\mathcal O_K$-module, and is
stable under division by $\pi$. Left and right multiplication by each
local unitary group preserve $W$.

Their $K$-linear span is the full algebra $D_K$. To check this, view
$\Res_{k/k_0}U$ as a group over $k_0$ and extend scalars from $K$
to an algebraic closure. Separability of $k/k_0$ and $L/k$ gives
independent pairs of matrix factors, on which a unitary
element has the form $(g,g^{-t})$ with $g\in\GL_n$. These pairs span
both matrix factors: varying a scalar multiple of $g$ separates the
weights $\lambda$ and $\lambda^{-1}$, and invertible matrices span
$M_n$. Independent factors give the whole product. The local points
of the smooth unitary groups are Zariski dense, so the same linear-span
assertion holds over $K$. It follows that $W$ is a two-sided ideal of
$D_K$.

Every simple factor occurs in this ideal. At a nonsplit place of $k$
the algebra over $L$ splits and the unitary group is isotropic because
$n\geq3$. Take an unbounded sequence in that local group and the
identity in every other component. For suitable integers
$r_\nu\to\infty$, multiplication by $\pi^{r_\nu}$ gives a subsequence
converging to a nonzero vector in that factor and to zero elsewhere.
For every fixed $r$, this sequence eventually belongs to $\pi^rC$;
closedness therefore puts its limit in $W$. At a split place, the
unitary group contains reciprocal central scalars. The same argument
with $(\pi^{-r},\pi^r)$, and then with the factors reversed, isolates
each of the two simple factors. Thus $W=D_K$, and $C=D_K$.

It remains to pass from all these local closures to an exact equality.
The $k_0$-linear span of $M$ is $D$, since a proper finite-dimensional
subspace would remain proper after completion. Choose a full
$R$-lattice $\Lambda\subset M$. For $x=\pi^{-j}\lambda$, with
$j\geq0$ and $\lambda\in\Lambda$, local density gives $a\in M$ with
$a-x\in\Lambda\otimes_R\mathcal O_K$. Choose a polynomial $q\in R$
which clears all denominators of $a$ away from $\pi$ and satisfies
$q\equiv1\pmod{\pi^j}$. The Chinese remainder theorem permits both
conditions. Then $qa-qx$ belongs to the local lattice at every finite
prime of $R$, hence to $\Lambda$, and $(q-1)x\in\Lambda$.
It follows that $x\in M$. Partial fractions now give every element of
$k_0\Lambda=D$ inside $M$, proving the lemma.
\end{proof}

\subsection{Separating the blocks by actual shears}

Skolem--Noether gives elements $w_j\in D^\times$ and a decomposition
\[
 D=\bigoplus_{j=1}^n E w_j,\qquad w_jd=t_jw_j.
\]
Thus $D\oplus D$ is the direct sum of $n$ two-dimensional $E$-blocks.
Let $H$ be the subgroup of the invariance group of $F_d$ generated by
$\diag(\lambda,\mu)$, with $\lambda,\mu\in E^1$, acting
simultaneously on every block, and by the simultaneous matrices
\begin{equation}\label{char:boosts}
 B_h(t_j)=\begin{pmatrix}1&t_j\bar h\\h&1\end{pmatrix}
 \qquad(h\in E^1).
\end{equation}
The phases act independently on the two $D$ coordinates; no independent
choice between blocks is assumed. These transformations are available:
changing the sign of the second
output in Equation~\ref{char:pair-move} gives $B_{-1}$, and diagonal
phase conjugation gives every $B_h$. Their determinants $1-t_j$ are
nonzero.

\begin{lemma}\label{char:supported-shears}
For every block $i$, $H$ contains a nonidentity upper elementary
unipotent supported only on block $i$, and a nonidentity lower
elementary unipotent supported only on block $i$.
\end{lemma}

\begin{proof}
We first construct a transformation triangular at a single block.
The first-column slope of $B_h(t)B_{hz}(t)$ is
$h(1+z)/(1+tz)$. At $t=t_j$, its norm from $E$ to $F$ is
\begin{equation}\label{char:slope-norm}
 \ell_j(z)=\frac{(1+z)^2}{(1+t_jz)(z+\bar t_j)}.
\end{equation}
This is a separable rational map of degree two. In characteristic two,
its derivative is nonzero because $t_j\bar t_j\ne1$; in other
characteristics separability is automatic. For a generic $z\in E^1$,
the other point of the fiber is
\[
 z'=\frac{(t_j-\bar t_j)z+1+t_j\bar t_j-2\bar t_j}
 {(1+t_j\bar t_j-2t_j)z+\bar t_j-t_j}.
\]
It is distinct from $z$ and lies in $E^1$: conjugate reciprocation
preserves the two roots and fixes $z$, so also fixes $z'$.
Writing $s_t(z)=(1+z)/(1+tz)$, choose
$h'=s_{t_j}(z)/s_{t_j}(z')\in E^1$ and set
\[
 g_j(t)=\bigl(B_{h'}(t)B_{h'z'}(t)\bigr)^{-1}B_1(t)B_z(t).
\]
Its lower-left numerator is
\[
 (1+tz')(1+z)-h'(1+z')(1+tz).
\]
This nonzero linear polynomial has its unique root at $t_j$; it is
nonzero as a polynomial because the two fractional slope functions
have distinct poles. Thus this entry is nonzero at every other
$t_\ell$ and every $\bar t_\ell$.

Matrices of the form
\[
 \begin{pmatrix}\alpha(t)&t\overline{\beta}(t)\\
 \beta(t)&\overline{\alpha}(t)\end{pmatrix},
\]
with bar acting on coefficients and fixing $t$, are closed under
multiplication and inversion. Therefore the upper-right entry of
$g_j(t)$ is nonzero at all $t_\ell$, including $t_j$.
All diagonal entries can also be kept nonzero. Indeed at $z=-1$ the
fiber interchange in Equation~\ref{char:slope-norm} fixes $-1$ and
has derivative $-1$ (equal to $1$ in characteristic two). Hence
$h'\to-1$, both boost products tend to $(1-t)I$, and $g_j(t)\to I$.
The excluded conditions are consequently proper algebraic conditions
on the rational norm-one curve, whose rational points are infinite.
We obtain an upper triangular nondiagonal matrix at $j$ and a matrix
with all four entries nonzero at each other block. Using $\bar t_j$
in the same construction gives the lower triangular analogue.

Fix $i$. For $j\ne i$, the subgroup generated by $g_j$ and diagonal
phases is triangular at $j$ and projectively Zariski dense at $i$.
For the second assertion, its closure contains the diagonal torus and
a matrix with four nonzero entries; a proper algebraic subgroup of
$\PGL_2$ containing that torus lies in a Borel subgroup or in the
normalizer of the torus, and neither contains such a matrix. The second derived
subgroup is trivial at $j$ but remains projectively dense at $i$,
by algebraic perfectness of $\PGL_2$.

Use the full normal kernels
\[
 K_j=\ker\bigl(H\longrightarrow\GL_2(E)
                    \text{ on block }j\bigr).
\]
Each $K_j$ has projectively dense image at $i$. Successive commutators
of these normal kernels are supported on the intersection of their
supports and still have projectively dense image at $i$, since
the commutator map is regular and its image generates the derived
algebraic group $\PGL_2$. More explicitly, successive groups
$N_{r+1}=[N_r,K_j]$, beginning with one $K_j$, remain in every kernel
already used because those kernels are normal. Consequently the image $P_i$ of
$\bigcap_{j\ne i}K_j$ is projectively dense and is normal in the full
image $H_i$ of $H$ at $i$.

We next produce actual unipotents in $P_i$. Commutators of the
triangular matrices at $i$ with diagonal phases give upper and lower
elementary unipotents in $H_i$ with nonzero coefficients. Conjugating
by phases and adding coefficients shows that both coefficient groups
contain scalar multiples of $\bbZ E^1$. We claim that they contain a
common nonzero ideal in a suitable ring of $S$-integers of $E$.
Choose a place of $F$ split in $E$ and a phase $z$ with a simple pole
above it, as in the proof of Lemma~\ref{char:additive-span}. Then
$t_0=z+z^{-1}\in F$ has a simple pole, so $E/\bbF_p(t_0)$ is finite
separable and $\bbZ E^1$ is a module over $R=\bbF_p[t_0]$.
The $k_0$-linear span of $E^1$ is $E$, where now
$k_0=\bbF_p(t_0)$. Indeed the group
$\Res_{F/k_0}(\ker(N_{E/F}:\Res_{E/F}\bbG_m\to\bbG_m))$
has rational points $E^1$. After extending scalars to an algebraic
closure of $k_0$, its coordinates are
independent pairs $(x_\ell,x_\ell^{-1})$. A linear relation
$\sum_\ell(a_\ell x_\ell+b_\ell x_\ell^{-1})=0$ forces all
coefficients to vanish, so these pairs span their ambient product.
Hilbert 90 gives Zariski density of its rational points. Thus
$\bbZ E^1$ contains a full $R$-lattice. If $\mathcal O$ is the
integral closure of $R$ in $E$, any two scalar multiples of such a
lattice contain a common nonzero ideal $I\subset\mathcal O$:
clear the finitely many denominators of generators of the finite
$R$-module $\mathcal O$ relative to both lattices.
There are distinct places $P,Q$ of $E$ above the pole of $t_0$.
The nonzero divisor $\deg(Q)P-\deg(P)Q$ has degree zero. Finiteness
of the degree-zero divisor class group over the finite constant field
gives a positive multiple equal to the divisor of a nonconstant
function. That function is an infinite-order unit of $\mathcal O$
\cite[Chapter~14]{Rosen2002}.

Choose $g\in P_i$ carrying $\infty$ to a finite point $r$ of
$\bbP^1(E)$. Fix $0\ne s\in I$ and choose an infinite-order unit
$u\in\mathcal O^\times$ sufficiently congruent to $1$ that
\[
 u=1+st\quad\text{for some }t\in I,
 \qquad (1-u^2)r\in I.
\]
A power of an infinite-order unit satisfies the finitely many required
congruences. With
$U(x)=\left(\begin{smallmatrix}1&x\\0&1\end{smallmatrix}\right)$
and
$L(x)=\left(\begin{smallmatrix}1&0\\x&1\end{smallmatrix}\right)$,
the identity
\[
 U(s)L(t)U(-s/u)L(-tu)=\diag(u,u^{-1})
\]
shows that $H_i$ contains the affine transformation
$h:x\mapsto u^2x+(1-u^2)r$. It fixes $r$ and $\infty$.
Normality puts $q=h^{-1}g^{-1}hg$ in $P_i$. This element fixes
$\infty$ and has affine slope $u^{-4}$: the derivatives of $h$ at
$\infty$ and $r$ are $u^{-2}$ and $u^2$, respectively.
Hence, for $0\ne x\in I$,
\[
 qU(x)q^{-1}U(-x)=U((u^{-4}-1)x)
\]
is a nonidentity upper unipotent in $P_i$. By definition of $P_i$
it lifts to a transformation supported only on block $i$.
Interchanging the two axes proves the lower assertion.
\end{proof}

\begin{proposition}\label{char:transitivity}
For $d\in E^\times$ satisfying
Equation~\ref{char:parameter-conditions}, the invariance group of $F_d$
is transitive on $D^2\setminus\{(0,0)\}$. Consequently, whenever
$b,bd\in\mathcal D$,
\begin{equation}\label{char:difference}
 f(bd)-f(b)=\chi(\bar d/d).
\end{equation}
\end{proposition}

\begin{proof}
If $c\ne0$, one of its $E$-block coordinates is nonzero. The pure
upper shear from Lemma~\ref{char:supported-shears} changes $a$ by
some fixed nonzero $y\in D$, leaving $c$ unchanged. Conjugating by
left multiplication of the first coordinate by $v\in U(k)$ replaces
this increment by $vy$. Products and inverses therefore add every
element of $(\bbZ U(k))y=D$, by Lemma~\ref{char:additive-span}.
The lower shears similarly permit arbitrary changes of $c$ when
$a\ne0$. These operations connect every nonzero pair to one with
both entries nonzero, and connect any two such pairs. Thus $F_d$ is
constant.

Take $e=(d-\bar d)^{-1}$, $c=0$ and $a=1$ in
Equation~\ref{char:sylvester}. The elements defining $f(e)$ and
$f(ed)$ are $1$ and $\bar d/d$, respectively. This evaluates the
constant. Any $b,bd\in\mathcal D$ yield a pair $(a,c)$ through their
defining equations; it is nonzero since the Sylvester operator has
zero kernel. This proves Equation~\ref{char:difference}.
\end{proof}

\subsection{From the difference identity to continuity}

Lemma~\ref{char:normalization} extends $f$ to every nonzero element of
$D$. For $b+b^*\ne0$, choose $s\in k^\times$ with $sb\in\mathcal D$
and put $f^\#(b)=f(sb)$; for skew $b$ put $f^\#(b)=0$.
This does not depend on the choice of $s$. If both $sb$ and $tb$ are
permitted, reduced norms of their defining Hermitian equations show
that $(t/s)^n$ is a norm from $L$. The quotient
$k^\times/N_{L/k}(L^\times)$ has exponent two, so oddness of $n$
implies that $t/s$ itself is a norm. Central congruence in
Equation~\ref{char:congruence} then gives the same value.
Consequently
\begin{equation}\label{char:sharp-invariance}
 f^\#(x^*bx)=f^\#(b),\qquad f^\#(sb)=f^\#(b)
 \quad(x\in D^\times,\ s\in k^\times).
\end{equation}

Let $m\geq1$ annihilate $B$, and let
$T=\ker(N_{E/L}:\Res_{E/L}\bbG_m\to\bbG_m)$.
Choose $d\in T(L)^m$ satisfying
Equation~\ref{char:parameter-conditions}. Such choices are Zariski
dense in a nonempty open subset of $\Res_{L/k}T$. Over a separable
closure of $k$, the two embeddings of $L$ give two independent
lists $(x_1,\ldots,x_n)$ and $(y_1,\ldots,y_n)$, each with product
one. These determinant-one eigenvalue lists can avoid all
stated equalities, and $n\geq3$ also allows their coordinatewise
products to be unequal. Hilbert 90 for the cyclic extension $E/L$
gives Zariski density of $T(L)$, and the power map is dominant even
when $p\mid m$.

For every such $d$ and every $b\in D^\times$,
\begin{equation}\label{char:step-invariance}
 f^\#(bd)=f^\#(b).
\end{equation}
Indeed normalize $b$ and $bd$ separately into $\mathcal D$.
Their ratio replaces $d$ by a nonzero $k$-multiple, preserving
Equation~\ref{char:parameter-conditions} and the ratio $\bar d/d$.
If $d=y^m$, then $\bar d/d=(\bar y/y)^m$ has zero character, so
Proposition~\ref{char:transitivity} applies. Congruence invariance
extends Equation~\ref{char:step-invariance} to every
$D^\times$-conjugate of $d$.

Let $J$ be the subgroup of $\SL_1(D)(L)$ generated by all these
conjugates. It is noncentral and normal. The established inner-type
Margulis--Platonov theorem over $L$ implies that $J$ contains every
rational point sufficiently close to $1$ at the places
\[
 A_D=\{w:\ D\otimes_L L_w\text{ is division}\}
\]
\cite[Sections~3.4--3.5]{PRsurvey}\cite{Rapinchuk2006}.
These are precisely the two places of $L$ over each place of $A$.
Since $f^\#$ is defined on all of $D^\times$, products of the
invariant steps give
\begin{equation}\label{char:inner-invariance}
 f^\#(br)=f^\#(b)\qquad(b\in D^\times,\ r\in J).
\end{equation}
We now use this identity to obtain the continuity of $\chi$ on $S(k)$.

\begin{proof}[Proof of Theorem~\ref{char:trivial}]
Choose $u\in U(k)$ with $\Nrd(u)\ne1$; a central phase suffices,
since the norm-one torus of $L/k$ has infinitely many rational points.
Division makes $1-u$ invertible. For $b=(1-u)^{-1}$,
\begin{equation}\label{char:final-b}
 b+b^*=1,\qquad b^*=-ub,\qquad f^\#(b)=\chi(u).
\end{equation}
If $u'\in U(k)$ has the same reduced norm as $u$, put
$b'=(1-u')^{-1}$. Oddness gives
$\overline{\Nrd(b)}=-\Nrd(u)\Nrd(b)$, and the same identity for $b'$,
so
\[
 q=\Nrd(b')/\Nrd(b)\in k^\times.
\]

For $u'$ sufficiently close to $u$ at $A$, choose $c\in D^\times$
with $\Nrd(c)=q^{(n-1)/2}$ and $c$ close to $1$ at $A_D$.
Existence without approximation follows from $H^1(L,\SL_1(D))=1$.
Locally the reduced norm is smooth, including when $p\mid n$:
its differential is the nonzero reduced trace. It therefore has
solutions near $1$ for norms near $1$. Weak approximation for
$\SL_1(D)$ adjusts the global solution to these local solutions.
Set $b''=q^{-1}c^*b'c$. Then
\[
 \Nrd(b'')=q^{-n}q^{n-1}\Nrd(b')=\Nrd(b),
 \qquad f^\#(b'')=f^\#(b').
\]
Thus $r=b^{-1}b''\in\SL_1(D)(L)$ is close to $1$ at $A_D$ and
belongs to $J$. Equation~\ref{char:inner-invariance} gives
$\chi(u')=\chi(u)$. Taking $u'=us$ proves continuity of
$\chi|_{S(k)}$ for the topology induced by $S_A$.

For clarity, this continuous finite-valued homomorphism extends to
$S_A$. Weak approximation makes $S(k)$ dense. If $K$ is its kernel,
then $\overline K\cap S(k)=K$ by continuity. The finite cosets of
$K$ cover $S(k)$; their closures therefore cover $S_A$. Thus
$\overline K$ is an open normal subgroup, and
$S(k)/K\simeq S_A/\overline K$, giving the extension.

At $v\in A$ write $U(k_v)=D_v^\times$ and
$S(k_v)=\SL_1(D_v)$. Approximate any two elements of $D_v^\times$
by rational unitary elements, prescribing the identity at the other
places of $A$. Their rational commutators have character zero, so
continuity kills every local commutator. By
Lemma~\ref{quo:local-abelianization}, these commutators generate a
dense subgroup of $\SL_1(D_v)$. The extension is therefore zero on
every factor of $S_A$, proving $\chi(S(k))=0$. If $A$ is empty,
the preceding continuity assertion already gives the result.
\end{proof}
